% concatenated from PrimeTilings.tex
\documentclass{amsart}
\usepackage{graphicx} % Required for inserting images
\usepackage{amsmath}  % must come before astheorem, not after 
\usepackage{amsthm}  
\usepackage{amsfonts}
\usepackage{amssymb}
\usepackage{comment}
\usepackage{enumerate}
\usepackage{xurl}
\usepackage[colorlinks=true,linkcolor=blue,citecolor=blue,urlcolor=blue]{hyperref}   % makes links work
\usepackage{placeins}  % \FloatBarrier


\def\vector#1#2#3{\left(
\begin{array}{l} 
#1\\ 
#2\\ 
#3 
\end{array}\right)}

\def\matrix#1#2#3#4#5#6#7#8#9{{
 \left( \begin{array}{ccc}
#1 & #2 & #3 \\
#4 & #5 & #6 \\
#7 & #8 & #9 
\end{array} \right)}}

\def\Erdos{Erd\H{o}s}

\usepackage{makecell}

\theoremstyle{plain}
\newtheorem{theorem}{Theorem}
%[section]
% Putting [theorem] in the next line means Lemmas are numbered
% in the same sequence with theorems.  [MB] 
\newtheorem{lemma}[theorem]{Lemma}
%[section]
\newtheorem{proposition}[theorem]{Proposition}
\newtheorem{corollary}[theorem]{Corollary}
%[section]

\theoremstyle{definition}
\newtheorem{definition}{Definition}[section]
%[section]
\newtheorem{example}{Example}[section]
\theoremstyle{remark} 
\newtheorem*{remark}{Remark}
\newcommand{\Q}{\mathbb{Q}} 


\def\Q{{\mathbb Q}}
\def\Z{{\mathbb Z}}
\def\ha{\frac{1}{2}}
\def\onehalf{\tfrac 1 2}
\def\cd{\cdot}

\title{Tiling a triangle into a prime number of congruent triangles}
\author[Beeson]{Michael Beeson}
\address{San Jos\'e State University (emeritus), and UCSC (research associate)}
\begin{document}


\begin{abstract}
We show that, apart from a few known exceptions,  if a triangle is 
cut into $N$ congruent triangles, then $N$ is not a prime number.  
The exceptions are cutting an isosceles triangle in half, cutting
an equilateral triangle in three, a 3-tiling of a 30-60-90 triangle,
and the long-known biquadratic tilings
of a certain right triangle, when $N$ is a sum of two squares.
\end{abstract}
\maketitle


\section{Introduction}

In our work on \Erdos\ problem~633 \cite{erdos-633, BLZ2026}, we had to show that triangles $T$ with 
incommensurable angles have non-square tilings.  We gave a formula for each 
possible shape of $T$, such that the
number of tiles $N$ and shape of the tile $(\alpha,\beta,\gamma)$ are 
determined by some rational parameters.  Any tiling by that tile must
satisfy the equation for some value(s) of the parameter(s), and at least
some (but not necessarily every) value of the parameters gives a tiling.
Here we use similar formulas for $N$, but instead of showing that $N$
is not square, we have to show that $N$ is not prime (except for $N$ equal to two or three).

There is a second \Erdos\ problem about triangle tiling, problem~634 \cite{erdos-634}, 
which asks for the characterization of all $N$ such that no triangle can be tiled into $N$
congruent triangles.  This paper shows that the prime numbers belonging to that set
are exactly those primes greater than three and congruent to 3 mod 4. 

In \cite{beeson-noseven}, we gave a direct and rather short proof that 
$N$ cannot be 7 or 11;  but the method could not reach 19.  

As a matter of notation,  we will always use $(a,b,c)$ for the side lengths
of the tile, and $\alpha,\beta,\gamma$ for the angles opposite those sides.

We divide our work into several  cases.  First we separate the cases when $T$ is
isosceles or equilateral.  After that we divide into cases
 according as $T$ has commensurable angles or not.
 The first two cases have been studied in prior works
 \cite{beeson-isosceles, beeson-isoscelesPrime,beeson-equilateral}, which will be 
 discussed below. 
 
 \begin{lemma} \label{lemma:RandT} 
 Let triangle $T$ be tiled by triangle $R$, where $T$ is not equilateral. Then $R$ has 
commensurable angles if and only if $T$ has commensurable angles.
\end{lemma}
\begin{proof} See the first paragraph of the proof of Theorem~5.3 in \cite{laczkovich1995}.
\end{proof}



The possible tilings with commensurable angles have been understood since the dawn
of the subject, and as we review below, in that case $N$ can be prime only when 
$N=2$ and $T$ is isosceles, or $N=3$ and $T$ is equilateral.  The fundamental 
observation of Laczkovich, elaborated in several papers, is that in the incommensurable case,
if $T$ is neither equilateral nor isosceles, then
we must have either 
\begin{itemize}
\item $3\alpha+2\beta = \pi$   \hskip3.1cm(Group 1), or 
\item $3\alpha + 3 \beta = \pi$,  i.e., $\gamma = 2\pi/3$  \hskip0.8cm(Group 2)
\end{itemize}
In each group,  there is a small finite number of possible shapes of $T$, since the 
angles of $T$ have to be positive integral combinations of $(\alpha,\beta,\gamma)$.
The argument is bulldozer-style at the top level:  we consider each possible shape
of tile and prove that $N$ cannot be prime.  Group~1 has already been dealt 
with in prior works (discussed below),  leaving Group~2 tilings to be dealt 
with in this paper.

Moreover, there is a close connection between ``incommensurable angles'' and 
``commensurable sides'':
\begin{theorem}\cite[Theorem 1.2]{beeson-zhang2026}
    \label{theorem:rationality}
 Let triangle $T$ be tiled by a tile $R$ such that
    \begin{itemize}
        \item $R$ is not similar to $T$;
        \item $R$ is not a right triangle;
        \item $R$ has incommensurable angles.
    \end{itemize}
    Then $R$ must have commensurable sides.
\end{theorem}

Once we have a tiling of some triangle $T$, and we know the sides of the tile $R$
are commensurable, then the sides of $T$ are commensurable too, since they are linear
integral combinations of the sides; after a suitable scaling the sides will be 
integers with no common factor.  That is how number theory enters this geometrical subject.

After the prior results are accounted for,  the only cases remaining to deal with in
this paper are when
$T$ has incommensurable angles and $\gamma = 2\pi/3$, and $T$ is not equilateral or isosceles.
 It is known that only four shapes
of $T$ are then possible.
Our plan is to prove that
\begin{itemize}
\item $(a,b,c)$ are commensurable, so we can assume they are integers.
\item The sides $(X,Y,Z)$ of $T$ are of the form $\lambda(X_0,Y_0,Z_0)$, where 
$(X_0,Y_0,Z_0)$ have no common factor.
\item $\lambda$ is an integer.
\item Then equating the area of $T$ to that of $N$ tiles, we have $N = \lambda^2 f(a,b)$ for
some function $f$, different for each shape.
\item $f$ is integer-valued (except in one case, which we will worry about later)
\item and finally, the value of $f(a,b)$ cannot be a prime.
\end{itemize}

Comparing this plan to the solution of \Erdos\ problem 633 in \cite{BLZ2026}:  the
problem there was to show that $N$ is not a square,  so instead of showing $f(a,b)$
is not a prime, we had to show $Z^2 = f(a,b)$ is not solvable.  Therefore the theory
of elliptic equations came into the proof.  In this paper, the number theory is
simpler, and the hard part comes from isosceles $T$.  In the \Erdos\ problem, 
isosceles $T$ is disposed of in one line, since every isosceles triangle can be
cut in two and two is not a square.  The proof that isosceles triangles
cannot be $N$-tiled for $N$ a prime greater than three has been given already,
in papers that will be cited here.  


\section{Prior results}
\begin{theorem} \label{theorem:triangletiling3}  Let $T$ be a triangle tiled into $R$
with  angles $(\alpha,\beta,\gamma)$ such that $$3\alpha + 2\beta = \pi,$$ and $T$
is not similar to $R$.
Then the number of tiles $N$ is not a prime.
\end{theorem}

\begin{proof} We first note that by Lemma~1 of \cite{beeson-triangletiling3},
the angles of $R$ are incommensurable, so we need not mention that in the 
hypothesis of this theorem.  The non-primality of $N$ is 
proved in \cite{beeson-triangletiling3}, by considering the 
possible shapes of $T$ one by one.   Specifically,
the possible shapes of triangles $T$ for a Group~1 tiling,  and the citation for 
the proof there is no prime tiling of that shape, are given in Table~\ref{table:Group1}. \qedhere
\end{proof}

\begin{table}[ht]
\caption{Group 1, $3\alpha + 2 \beta = \pi$}
\begin{center}
\begin{tabular}{c l}
Angles of $T$ &  citation for $N$ not prime\\
\hline
$(2\alpha,\alpha,2\beta)$  & \cite{beeson-triangletiling3}, Theorem~12 \\
$(2\alpha, \beta, \alpha + \beta)$ &  \cite{beeson-triangletiling3}, Theorem~8\\
$(\beta,\beta,3\alpha)$  & \cite{beeson-triangletiling3}, Theorem~15\\
$(\alpha+\beta,\alpha+\beta,\alpha)$  & \cite{beeson-triangletiling3}, Theorem~18\\
$(\alpha,\alpha,\alpha+2\beta)$  & \cite{beeson-triangletiling3}, Theorem~20
\end{tabular}
\end{center}
\label{table:Group1}
\end{table}%




\begin{theorem} \label{theorem:equilateral}  Let $T$ be an equilateral triangle $N$-tiled, for $N \neq 3$.
Then the number of tiles $N$ is not a prime.
\end{theorem}

\begin{proof} This is proved in \cite{beeson-equilateral}.  Laczkovich already
showed that $\gamma$ must be $\pi/3$ or $2\pi/3$ \cite{laczkovich1995, laczkovich2012}, especially Theorem~3.3 of \cite{laczkovich2012}.
Specifically, Laczkovich showed that one of following holds:

\medskip

(i) $(\pi/3,\pi/3,\pi/3)$  (equilateral)

(ii) $(\pi/6,\pi/6,2\pi/3)$
\smallskip

(iii) $(\pi/6, \pi/2, \pi/3)$
\smallskip

(iv) $(\alpha, \beta, \pi/3)$ with $\alpha$ not a rational multiple of $\pi$
\smallskip

(v)  $(\alpha,\beta,2\pi/3)$ with $\alpha$ not a rational multiple of $\pi$
\smallskip

The angles are commensurable in the first three cases, and incommensurable 
in the last two. 
The possible $N$ in the first three cases are known.  In case (i),
$N$ must be a square, and any square corresponds to a tiling.  (That 
is true for any triangle $ABC$ and a tile similar to $ABC$.)  In cases (ii) and (iii),
$N$ must have the form $3n^2$ or $6n^2$ respectively.
None of these can be prime except for case (ii), when $N=3$ is possible.

Explicit citations for the nonprimality of $N$ in these cases
are given in Table~\ref{table:equilateral}.
\begin{table}[ht]
\caption{Equilateral $T$}
\begin{center}
\begin{tabular}{r l}
Angles  &  citation for $N$ not prime\\
%\hline
$\alpha = \beta = \pi/3$  & \cite{snover1991} or \cite{soifer} \\
$\alpha = \beta = \pi/6$  & \cite{laczkovich2012}, Theorem~3.1 \\
$(\alpha,\beta,\gamma) =(\pi/6, \pi/2, \pi/3)$   & \cite{laczkovich2012}, Theorem~3.1  \\
$\gamma =\pi/3, \alpha \not\in \Q$ & \cite{beeson-triangletiling3}, Theorem~3\\
$\gamma = 2\pi/3, \alpha \not\in \Q$  & \cite{beeson-triangletiling3}, Theorem~6\
\end{tabular}
\end{center}
\label{table:equilateral}
\end{table}%
\end{proof}


\begin{theorem} \label{theorem:isosceles}  Let $T$ be an isosceles (but not equilateral) triangle $N$-tiled
using a tile $R$ with angles $(\alpha,\beta,\gamma)$. Suppose $T$ is not similar to $R$ and $N \neq 2$. Then
\begin{enumerate}[{\rm (i)}]
\item $\gamma = 2\alpha$, or $\gamma = \pi/2$, or $\gamma = 2\pi/3$, or $3\alpha + 2\beta = \pi$.
\item Unless $\gamma = \pi/2$, the angles of $R$ are incommensurable.
\item $N$ is not prime.
\end{enumerate}
\end{theorem}

\begin{proof}
Parts~(i) and (ii) are proved in \cite{beeson-equilateral}, resting on \cite{laczkovich1995}.
Turning to part~(iii), we note that the base angles of $T$ are $\alpha$ or $\beta$.  Except
in case $3\alpha+2\beta = \pi$, we may assume without loss of generality that the base
angles are $\alpha$.  Citations for the non-primality of $N$ in each case are 
given in Table~\ref{table:isosceles}. 
\end{proof}
\begin{table}[ht]
\caption{$T$ isosceles but not equilateral}
\begin{center}
\begin{tabular}{c l}
Angles &  citation for $N$ not prime\\
\hline
$\gamma = 2 \alpha$  & \cite{beeson-isosceles}, Theorem~11.7  \\
$3\alpha +2\beta = \pi$ &  see Theorem~\ref{theorem:triangletiling3} above\\
$\gamma = \pi/2$   &  \cite{beeson-isosceles}, Corollary~7.9 \\
$\gamma = 2\pi/3$  &  \cite{beeson-isoscelesPrime}, Theorem~11
\end{tabular}
\end{center}
\label{table:isosceles}
\end{table}%


A {\em reptiling} is a tiling in which the tile is similar to the tiled triangle.
\begin{theorem}[\cite{laczkovich1995,laczkovich2012}]
   \label{theorem:l1995-5.3} Suppose $T$ has commensurable angles and tiles into $R$.  If $T$ is not an isosceles (or equilateral) triangle, then the tiling is a reptiling.
\end{theorem}

\begin{proof} Following \cite{laczkovich1995}, let $c(T)$ be the 
number of distinct non-similar triangles $R$ into which $T$ can be tiled.  Then $c(T) = 1$ 
means that $T$
can only be tiled by $R$ similar to $T$ (that is, a reptiling). 
Theorem~5.3 of \cite{laczkovich1995} says that, under the assumptions that $T$ has
commensurable angles and is not isosceles, or is the isosceles right triangle,
then $c(T) = 1$, and if $T$ is isosceles but neither right nor equilateral, then $c(T) = 2$.
Hence, under the assumptions of our theorem, $c(T) = 1$. 
\end{proof}

The main result related to reptiling is:
\begin{theorem}[\cite{snover1991}]
    \label{theorem:reptile} Suppose $T$ can be reptiled into $N$ tiles. Then $N$ must be one (or more) of:
    \begin{itemize}
        \item $N = M^2$, which is possible for any triangle $T$ and any natural number $M$;
        \item $N = M^2 + K^2$, in case where $N$ is not a square, $T$ must be a right triangle with legs having ratio $M/K$;
        \item $N = 3M^2$, in which case $T$ must be the $(\pi/2, \pi/3, \pi/6)$ right triangle.
    \end{itemize}
\end{theorem}

Consequently:
\begin{theorem} \label{theorem:reptilenotprime}  
Suppose $T$ can be reptiled into $N$ tiles and $N \not\in\{2, 3\}$ is prime. 
Then $N \equiv 1 \pmod 4$, and 
$T$ is a right triangle with legs in ratio $M/K$.
\end{theorem}

\begin{remark} A 30-60-90 triangle can be 3-reptiled, and a right isosceles triangle is 2-reptiled by its altitude.
See Figure~\ref{figure:tiling13}.   
\end{remark}

\begin{proof}  Assume $N$ is a prime number.
Looking at the three cases of Theorem~\ref{theorem:reptile}, $N$ is a square, or three times a square, or 
a sum of two squares.  In the first two cases, $N$ is not prime, except when $N=3$.
  Therefore, $N$  is a sum
of two squares $M^2 + K^2$, and $T$ is a right triangle with legs in ratio $M/K$.  Therefore
$N \equiv 1 \pmod 4$, since $N \neq 2$ by hypothesis.
\end{proof}

\begin{figure}[ht]
\includegraphics[width=0.2\textwidth]{tiling2}
\hskip0.5cm
\includegraphics[width=0.3\textwidth]{tiling3}
\hskip0.5cm
\includegraphics[width=0.3\textwidth]{tiling13}
\caption{Prime tilings}
\label{figure:tiling13}
\end{figure}


Note that if $T$ tiles into congruent copies of $R$ and $T$ has incommensurable angles, then $R$ must also have incommensurable angles; otherwise, each angle of $T$ would be equal to a sum of rational multiples of $\pi$ from the tiling, so $T$ would have commensurable angles. This allows us to use the following corollary of another result from \cite{laczkovich1995}: 

\begin{theorem}[Corollary of Theorem 4.1, \cite{laczkovich1995}]
\label{theorem:l1995-4.1}
Let a non-isosceles triangle $T$ be tiled into $R$ where $R$ has incommensurable angles $(\alpha, \beta, \gamma)$. Then one of the following must hold:
\begin{enumerate}[{\rm (i)}]
    \item the tiling is a reptiling;
    \item (Group 1): $3\alpha + 2\beta = \pi$ and $T$ has angles $(\alpha, 2\alpha, 2\beta)$ or $(2\alpha, \beta, \alpha + \beta)$; or
    \item (Group 2): $3\alpha + 3\beta = \pi$ (equivalently, $\gamma = 2\pi/3$) and $T$ has angles 
    $$(\alpha, 2\alpha, 3\beta), (\alpha, 2\beta, 2\alpha + \beta), (\alpha, \alpha+\beta, \alpha+2\beta), \text{or }  (2\alpha, 2\beta, \alpha+\beta).$$
\end{enumerate}
\end{theorem}


\begin{proof}
See \cite{laczkovich1995}, Theorem~4.1.
That theorem has more cases in the conclusion than
shown here, but does not have our hypothesis that 
$T$ is not isosceles. 
By adding in the requirement that $T$ is not isosceles, we reduce the number of cases of the original result. 
\end{proof}


\section{Some lemmas about a tile with a $2\pi/3$ angle}
Here we present some lemmas that are only about the tile,
and not about the shape of the tiled triangle.


\begin{lemma} \label{lemma:claude}
Let triangle $R$ have angles $(\alpha,\beta,2\pi/3)$ and sides $(a,b,c)$ opposite
those angles, respectively. Then
\begin{eqnarray*}
\cos \alpha &=& \frac{a+2b} {2c} \\
\sin \alpha &=& \frac {a \sqrt 3}{2c}\\
\sin 2 \alpha &=& \frac {a \sqrt 3 (a+2b)}{2c^2}\\
\sin 3 \alpha 
&=&  \frac {a \sqrt 3}{2c} \cdot\frac {3b(a+b)}{c^2}
\end{eqnarray*}
\end{lemma}
 
\begin{proof}
By the law of sines, $\sin \alpha = (a/c) \sin (2\pi/3) = \frac {a \sqrt 3} {2c}$, as claimed.

To evaluate $\cos \alpha$: 
\begin{eqnarray*}
\cos^2 \alpha &=& 1- \sin^2 \alpha = 1- \frac {3a^2}{4c^2} = \frac{4c^2-3a^2}{4c^2} = \frac{4(a^2+b^2+ab)-3a^2}{4c^2} \\
\cos^2 \alpha  &=& \frac{a^2 + 4b^2+4ab}{4c^2} = \frac{(a+2b)^2}{4c^2} \\
\cos \alpha &=& \frac {a+2b} {2c} \mbox{\qquad as claimed}
\end{eqnarray*}
Then we have
\begin{eqnarray*}
\sin 2\alpha &=& 2 \sin \alpha \cos \alpha = 2\frac {a \sqrt 3} {2c} \cdot\frac{a+2b} {2c} = \frac {a \sqrt 3 (a+2b)}{2c^2}\\
\end{eqnarray*}
and finally 
\begin{eqnarray*}
\sin 3 \alpha &=& \sin \alpha (3 - 4 \sin^2 \alpha) \\
&=&  \frac {a \sqrt 3}{2c} \left(3 - \frac {3a^2}{c^2}\right)  = \frac {a \sqrt 3}{2c} \left( \frac{3c^2 - 3a^2}{c^2}\right)\\
&=&  \frac {a \sqrt 3}{2c} \cdot \frac {3b(a+b)}{c^2} \qedhere
\end{eqnarray*}
\end{proof}

\begin{lemma}\label{lemma:beta}
Let triangle $R$ have angles $(\alpha,\beta,2\pi/3)$ and sides $(a,b,c)$ opposite
those angles, respectively. Then
\begin{eqnarray*}
\sin \beta &=& \frac{b \sqrt 3}{2c}\\
\cos \beta &=& \frac {2a+b}{2c}\\
\sin 2 \beta &=& \frac {b \sqrt 3 (2a+b)}{2c^2}\\
\sin (2 \alpha +\beta) &=& \frac{\sqrt 3 (a+b)}{2c} \\
\sin (\alpha + 2 \beta)&=&\frac{\sqrt 3 (a+b)}{2c} 
\end{eqnarray*}
\end{lemma}

\begin{proof}
The first three formulas can be obtained by switching the labels $a$ and $b$, and $\alpha$ and $\beta$, in Lemma~\ref{lemma:claude}, 
which is legal since the requirement $3\alpha + 3\beta = \pi$ is invariant under that switch.  Then,
$$ \sin(2\alpha + \beta) = \sin(2\pi/3 -\alpha)= \frac{\sqrt 3} 2 \cos \alpha + \frac 1 2 \sin \alpha  = \sqrt 3 \frac {a+b}{2c}. $$
Finally, $\sin(\alpha+2\beta) = \sin(2\alpha+\beta)$, since the sum of those two angles is $\pi$.
\end{proof}

\begin{lemma} \label{lemma:pairwisecoprimality} Suppose $c^2 = a^2 + b^2 + ab$ and 
$(a,b,c)$ have no common divisor.  Then $(a,b,c)$ are pairwise relatively prime.
\end{lemma}
\begin{proof}  Suppose $p$ is prime and $p \mid a$ and $p \mid b$.  Then $p \mid a^2 + b^2 + ab$, so
$p \mid c$, so $p$ divides all three of $(a,b,c)$, contradicting the hypothesis.
  Suppose instead that $p \mid a$ and $p \mid c$.  Then $p \mid c^2 -a^2 - ab = b^2$,
so $p \mid b$, again contradicting the hypothesis.  Similarly with $a$ and $b$ switched. 
\end{proof}

\begin{lemma} \label{lemma:BLZ30} Suppose $c^2 = a^2 + b^2 + ab$ and 
$(a,b,c)$ have no common divisor. 
$$ \gcd(c^2,c(a+2b),3b(a+b)) = 1.$$
\end{lemma}

\begin{proof}  This is proved in \cite{BLZ2026}, Lemma~30.
\end{proof}

\begin{lemma} \label{lemma:claude3}
Suppose triangle $T$ is tiled by a triangle $R$ with integer sides.
 Let the side lengths 
$(X,Y,Z)$ of $T$ be proportional to coprime integers $(X_0, Y_0, Z_0)$.
Then $(X,Y,Z) = \lambda (X_0,Y_0,Z_0)$ for some integer $\lambda$.
\end{lemma}

\begin{proof} Since there is a tiling, each side length of $T$ is a 
linear integral combination of $(a,b,c)$, which are integers by 
hypothesis.  Hence $(X,Y,Z)$ are integers.  By hypothesis,
$$ (X,Y,Z) = \lambda (X_0,Y_0,Z_0)$$
for some $\lambda$.  Then $\lambda$ is rational, since $X$ and $X_0$
are integers.  Let $p,q$ be coprime positive integers with $\lambda = p/q$.
We have 
$$ q (X,Y,Z) = p (X_0,Y_0,Z_0).$$
Since $p$ and $q$ are coprime, $q$ divides each of $(X_0$, $Y_0$, $Z_0)$.  But 
by hypothesis, these three have no common factor.  Hence $q = 1$ and $\lambda = p$
is an integer. 
\end{proof}

\begin{lemma} \label{lemma:claude4} Let  $(a,b,c)$ be positive integers
satisfying $c^2 = a^2 + ab + b^2$.  Then $a+b$ is composite.
\end{lemma}

\begin{proof}  Suppose $p := a+b$ is prime.  Then 
$c^2 = p^2 - ab$, and  we have
\begin{eqnarray*}
\tfrac 3 4 p^2 \le c^2 < p^2   &&  \mbox{\qquad since 
$0< p^2-c^2 = ab \le \frac {(a+b)^2} 4 = p^2/4$}\\
p/2 < c < p &&   \mbox{\qquad since $c^2 = p^2-ab$}\\
b \equiv -a \pmod p   &&\mbox{\qquad since $p:= a+b$}\\
c^2 \equiv p^2-ab \equiv a^2 \pmod p && \mbox{\qquad since $b \equiv -a$} \\
c^2 - a^2 \equiv 0 \equiv (c-a)(c+a)  && \mbox{\qquad since $c^2 \equiv a^2$}\\
p | (c-a)(c+a)  && \\
|c-a| < p  &&  \mbox{\qquad since $0 < a < p$ and $0 < c < p$} \\
c\neq a && \mbox{\qquad since if $c=a$ then $b(a+b)=0$, which is false}\\
p \nmid (c-a)  && \mbox{\qquad by the preceding lines}\\
p | (c+a)     && \mbox{\qquad since $ p | (c-a)(c+a)$ but $p \nmid (c-a)$}\\
c+a = p = a+b   && \mbox{\qquad since $0 < c+a < 2p$}\\
c = b    && \\
a(a+b) = 0  && \mbox{\qquad since $c^2 = a^2 + ab + b^2$ and $c=b$} 
\end{eqnarray*}
But that is impossible, since $a$ and $b$ are positive.
\end{proof}

\begin{lemma} \label{lemma:not3divc} Let $(a,b,c)$ have no common factor
and satisfy $c^2 = a^2 + ab + b^2$.  Then $3 \nmid c$.
\end{lemma}

\begin{proof}
By Lemma~\ref{lemma:pairwisecoprimality}, $\gcd(b,c) = 1$. 
Suppose $3 | c$. Then $9 | c^2 = a^2 + ab+b^2$.   Then $a^2 + ab + b^2 \equiv 0 \pmod 3$.
We will show $ a \equiv b \pmod 3$.  If $a \equiv 0$, then $0 \equiv a^2 + ab + b^2 \equiv b^2$, so $b \equiv 0 \equiv a$.
If $a \equiv 1$, then $0 \equiv a^2 + ab + b^2 = 1+ b + b^2$.  Then $b \equiv 1 \equiv a$.  If $a=2$ then 
$0 \equiv 1 + 2b + b^2$, which implies $b \equiv 2 \equiv a$.  Therefore $a \equiv b$, as claimed.   Therefore
there exists an integer $t$ such that $b = a + 3 t$.   Then 
$$ a^2 + ab + b^2 \equiv 3 b^2 + 9 bt + 9 t^2  \equiv 3b^2 \pmod 9.$$
But $3b^2$ is nonzero (mod 9), since $3 | c$ and $\gcd(b,c) = 1$. Therefore $9 \nmid c^2$, contradiction.
\end{proof}



\section{The case $(\alpha,2\alpha,3\beta)$}

\begin{lemma}\label{lemma:rationality3}
Suppose triangle $T$ with incommensurable angles is tiled using triangle $R$ with angles $(\alpha, \beta, 2\pi/3)$.
Suppose $T$ is not similar to $R$.
Then the side lengths $(a,b,c)$ of $R$ are commensurable.
\end{lemma}

\begin{proof}
$R$ is not a right triangle, since $\gamma = 2\pi/3$ and $\alpha, \beta < \pi/3$.
$T$ is not equilateral, since its angles are incommensurable.
$R$ has incommensurable angles, by Lemma \ref{lemma:RandT}.
$R$ is not similar to $T$, by hypothesis.  Then by Theorem~\ref{theorem:rationality},
$(a,b,c)$ are commensurable.
\end{proof}

\begin{theorem} \label{theorem:case2}  Let triangle $T$ have incommensurable angles 
$(\alpha, 2\alpha, 3\beta)$, and suppose $T$ is $N$-tiled  by $(\alpha,\beta,2\pi/3)$.
Then $N$ is not a prime.  Moreover,
$$ N =  3\lambda^2 (a+2b)  (a+b),$$
where $(a,b,c)$ are the sides of the tile with no common factor and $\lambda \in \Z$.
\end{theorem}


\begin{proof} $T$ has no $2\pi/3$ angle, since if $3\beta = 2\pi/3$ then $\beta = 2\pi/9$,
violating the incommensurable angles hypothesis; and if $2\alpha = 2\pi/3$ then $\alpha = \pi/3$,
also violating the incommensurable angles hypothesis; and if $\alpha = 2\pi/3$ then $2\alpha > \pi$
could not be an angle of $T$.
Since $T$ has no $2\pi/3$ angle, it is not similar to the tile.
By Lemma~\ref{lemma:rationality3}, $(a,b,c)$ are commensurable.
Without loss of generality, we may assume they are integers with no common factor.
Then the side lengths $(X,Y,Z)$ are also integers, since they are integral linear
combinations of $(a,b,c)$. 
Let the sides of $T$ next to the $3\beta$ angle be $X$ and $Y$. 
By the law of sines,
$$ X: {\sin \alpha} =   Y :{\sin 2 \alpha} = Z :{\sin 3 \beta}.$$
Since $\pi = 3 \alpha + 3\beta$, we have $\sin 3\beta = \sin 3 \alpha$.   Therefore
$$ X: {\sin \alpha} =   Y :{\sin 2 \alpha} = Z : {\sin 3 \alpha}.$$
By Lemma~\ref{lemma:claude},
$$ X : \frac {a \sqrt 3}{2c} = Y:\frac {a \sqrt 3 (a+2b)}{2c^2} = Z: \frac {a \sqrt 3}{2c} \cdot \frac{3b(a+b)}{c^2}.$$
Dividing by $a \sqrt{3}/2$ and multiplying by $c^3$ we have
$$ X : c^2 = Y: c(a+2b) = Z:  {3b(a+b)}.$$
That is, for some real number $\lambda$, we have 
\begin{equation}
 (X,Y,Z) = \lambda(c^2,c(a+2b),3b(a+b)). \label{eq:421}
 \end{equation}
Since $X$ and $c$ are positive integers, $\lambda$ is a positive rational.
By Lemma~\ref{lemma:BLZ30}, $\gcd(c^2,c(a+2b),3b(a+b)) = 1$. 
By Lemma~\ref{lemma:claude3}, $\lambda$ is a positive integer.
(This is the key step of the proof.)

 The area equation expresses that the area of $T$ is $N$ times the area of the tile:
\begin{equation*}
XY \sin 3\beta = N ab \sin \gamma \ = \ N ab \frac {\sqrt 3} 2  
\end{equation*}

Substitute these values for $X$ and $Y$ from (\ref{eq:421}) into the area equation.
We obtain
\begin{eqnarray*}
XY \sin 3\beta &=& \ N ab \frac {\sqrt 3} 2  \\
\lambda^2 c^3(a+2b)\sin 3 \beta &=& Nab \frac{\sqrt 3} 2
\end{eqnarray*}
Since $3\alpha+3\beta =\pi$, we have $\sin 3 \beta = \sin 3 \alpha$.  Therefore
\begin{eqnarray*}
\lambda^2 c^3(a+2b)\sin 3 \alpha &=& Nab \frac{\sqrt 3} 2
\end{eqnarray*}
Substituting for $\sin 3 \alpha$ from Lemma~\ref{lemma:claude}, we have
\begin{eqnarray*}
\lambda^2 c^3(a+2b) \cdot \frac {a \sqrt 3}{2c} \cdot \frac {3b(a+b)}{c^2} &=& Nab \frac{\sqrt 3} 2 \\
 3\lambda^2 (a+2b)  (a+b) &=& N
\end{eqnarray*}
Since $\lambda$ is an integer and $a$ and $b$ are positive integers, we have here a factorization of $N$.
Hence $N$ is not prime.
\end{proof}



\section{The case $(\alpha,2\beta,2\alpha+\beta)$}


\begin{theorem} \label{theorem:case3}  Let triangle $T$ have incommensurable angles 
$(\alpha, 2\beta, 2\alpha + \beta)$, and suppose $T$ is $N$-tiled by $(\alpha,\beta,2\pi/3)$.
Then $N$ is not a prime.  Moreover,
$$ N = \lambda^2 (2a+b) (a+b),$$
where $(a,b,c)$ are the sides of the tile with no common factor, and $\lambda \in \Z$.
\end{theorem}

\begin{proof}  Since $\alpha, \beta < \pi/3$,  neither $\alpha$ nor $2\beta$ is equal to $2\pi/3$.
If $2\alpha + \beta = 2\pi/3$ then $\alpha = \pi/3$, since $\alpha + \beta = \pi/3$; but that
violates the incommensurable angles hypothesis.  Hence $T$ has no $2\pi/3$ angle. 
Therefore $T$ is not similar to the tile.
By Lemma~\ref{lemma:rationality3}, $(a,b,c)$ are commensurable.
Without loss of generality, we may assume they are integers with no common factor.
Then the side lengths $(X,Y,Z)$ are also integers, since they are integral linear
combinations of $(a,b,c)$. 
Let the sides of $T$ next to the $2\alpha + \beta$ angle be $X$ and $Y$.  
We will show that $(X,Y,Z)$  are proportional to a triple with
no common factor, apply Lemma~\ref{lemma:claude3} to conclude
that the proportionality factor is an integer, 
and then use the area equation to finish.   Here are the details:

By the law of sines,
$$ X: {\sin \alpha} =   Y :{\sin 2 \beta} = Z :{\sin( 2\alpha + \beta)}.$$
Using the formulas from Lemmas~\ref{lemma:claude} and~\ref{lemma:beta}, we have 

$$ X: \frac {a \sqrt 3}{2c} = Y:\frac {b \sqrt 3 (2a+b)}{2c^2} = Z: \frac{\sqrt 3 (a+b)}{2c}.$$
$$ X: ac = Y : b(2a+b) = Z: (a+b)c.$$
Then for some real $\lambda$ we have
\begin{equation}
 (X,Y,Z) = \lambda (ac, b(2a+b), (a+b)c). \label{eq:518}
 \end{equation}
We will prove $(ac,b(2a+b),(a+b)c)$ have no common factor. 
We have $c^2 = a^2 + b^2 + ab$; by Lemma~\ref{lemma:pairwisecoprimality},
we have $\gcd(a,b) = 1$. 
Suppose the prime $p$ divides all three of $(ac,b(2a+b),(a+b)c)$.
  Then $p | ac$, so 
$p | a$ or $p | c$.  Then $p \nmid b$.  But $p | b(2a+b)$,
so $p | (2a+b)$.  If $p | a$ then $p | b$; contradicting $\gcd(a,b) = 1$. 
Therefore $p \nmid a$. Therefore $p | c$, since $p | a$ or $p | c$. 
Modulo $p$ we have
\begin{eqnarray*}
b \equiv -2a   && \mbox{\qquad since $p | (2a+b)$}\\
0 \equiv c^2 \equiv a^2 + b^2 + ab && \mbox{\qquad since $p | c$}\\
0  \equiv 3a^2 \pmod p  && \mbox{\qquad since $b \equiv -2a$}\\
p = 3                && \mbox{\qquad since $p \nmid a$}\\
\end{eqnarray*}
But that contradicts Lemma~\ref{lemma:not3divc}.
Therefore $(ac,b(2a+b),(a+b)c)$ have no common factor. 
Then by Lemma~\ref{lemma:claude3}, $\lambda$ is a positive integer in (\ref{eq:518}).

The area equation for this shape of triangle is 
$$ XY \sin (2\alpha+\beta) \ = \ N ab \frac {\sqrt 3} 2.$$
Substituting the values for $X$ and $Y$ from (\ref{eq:518}) into the area equation,
we obtain
\begin{eqnarray*}
\lambda^2 abc(2a+b) \sin(2\alpha+\beta) &=& \ N ab \frac {\sqrt 3} 2  \\
\lambda^2 c(2a+b) \sin(2\alpha+\beta)  &=& N \frac {\sqrt 3} 2  
\end{eqnarray*}
Putting in the value of $\sin(2\alpha+\beta)$ from Lemma~\ref{lemma:beta}, we have
\begin{eqnarray*}
\lambda^2 c(2a+b) \sqrt 3 \frac {a+b}{2c}  &=& N \frac {\sqrt 3} 2  \\
\lambda^2 (2a+b) (a+b)  &=& N 
\end{eqnarray*}
Since $\lambda$ is an integer, and $a,b,c$ are positive, we have a factorization of $N$.
Hence $N$ is not prime.
\end{proof}




\section{The case $(\alpha, \alpha+\beta, \alpha+2\beta)$}
\begin{theorem} \label{theorem:case4}  Let triangle $T$ have incommensurable angles 
$(\alpha, \alpha + \beta, \alpha + 2\beta)$, and suppose $T$ is $N$-tiled by $(\alpha,\beta,2\pi/3)$.
Then $N$ is not a prime.  Moreover, 
$$ N = \frac {\lambda^2 (a+b)} b,$$
where $(a,b,c)$ are the sides of the tile with no common factor and $\lambda \in \Z$.
\end{theorem}

\begin{proof}  
 Since $\alpha, \beta < \pi/3$,  neither $\alpha$ nor $\alpha + \beta$ is equal to $2\pi/3$.
If $\alpha + 2\beta = 2\pi/3$ then $\beta = \pi/3$, since $\alpha + \beta = \pi/3$; but that
violates the incommensurable angles hypothesis.  Hence $T$ has no $2\pi/3$ angle. 
Therefore $T$ is not similar to the tile.
By Lemma~\ref{lemma:rationality3}, $(a,b,c)$ are commensurable.
Without loss of generality, we may assume they are integers with no common factor.
Then the side lengths $(X,Y,Z)$ are also integers, since they are integral linear
combinations of $(a,b,c)$. 
Let the sides of $T$ next to the $\alpha + 2\beta$ angle be $X$ and $Y$.  
Again we will show that $(X,Y,Z)$ are multiples of a triple with no common factor, and apply Lemma~\ref{lemma:claude3} 
and then the area equation.   Here are the details:

By the law of sines,
$$ X: {\sin \alpha} =   Y :{\sin(\alpha + \beta)} = Z :{\sin( \alpha + 2\beta)}.$$
Using the formulas from Lemmas~\ref{lemma:claude} and~\ref{lemma:beta},  as well as $\alpha + \beta = \pi/3$, we have

$$ X: \frac {a \sqrt 3}{2c} = Y:\frac { \sqrt 3 }2 = Z: \frac{\sqrt 3 (a+b)}{2c} .$$
$$ X: a = Y : c = Z: (a+b).$$
Then for some real $\lambda$ we have
\begin{equation}
 (X,Y,Z) = \lambda (a, c, a+b). \label{eq:575}
 \end{equation}
Now $(a, c, a+b)$ have no common factor, since if prime $p$ divides $a$, $c$, and $a+b$, then it also divides $b$,
 contradiction.  
Then by Lemma~\ref{lemma:claude3}, $\lambda$ is a positive integer in (\ref{eq:575}).

The area equation for this shape of triangle is 
$$ XY \sin (\alpha+2 \beta) \ = \ N ab \frac {\sqrt 3} 2.$$
Substituting the values for $X$ and $Y$ from (\ref{eq:575}) into the area equation,
we obtain
\begin{eqnarray*}
\lambda^2 ac \sin(\alpha+2 \beta) &=& \ N ab \frac {\sqrt 3} 2  \\
\lambda^2 c \sin(\alpha+2\beta)  &=& Nb \frac {\sqrt 3} 2  
\end{eqnarray*}
Putting in the value of $\sin(\alpha+2\beta)$ from Lemma~\ref{lemma:beta}, we have
\begin{eqnarray*}
\lambda^2 c \sqrt 3 \frac {a+b}{2c}  &=& Nb\frac {\sqrt 3} 2  \\
N &=& \frac {\lambda^2 (a+b)} b
\end{eqnarray*}
We have $\gcd(a+b,b) = \gcd(a,b) = 1$.  Hence $\lambda^2/b$ is an integer. 
 By Lemma~\ref{lemma:claude4}, $a+b$ is composite.  
Hence $N$ is not prime.
\end{proof}


\section{The case $(2\alpha, 2\beta, \alpha+\beta)$}
\begin{theorem} \label{theorem:case5}  Let triangle $T$ have incommensurable angles 
$(2\alpha, 2\beta, \alpha + \beta)$, and suppose $T$ is $N$-tiled by $(\alpha,\beta,2\pi/3)$.
Then $N$ is not a prime.  Moreover,
$$N = \lambda^2 (a+2b) (2a+b),$$
where $(a,b,c)$ are the sides of the tile with no common factor and $\lambda \in \Z$. 
\end{theorem}

\begin{remark} This is Yan's Conjecture~5.1 in \cite{zhang2026}. He proved
this formula under the assumption $a \not\equiv b \pmod 3$.  Lemma~\ref{lemma:not3divc}
renders that assumption unnecessary; but here we give an independent proof. 
\end{remark}

\begin{proof} 
 Since $\alpha, \beta < \pi/3$,  neither $2\alpha$ nor $2\beta$ nor $\alpha+\beta$ is equal to $2\pi/3$. 
Therefore $T$ is not similar to the tile.
By Lemma~\ref{lemma:rationality3}, $(a,b,c)$ are commensurable.
Without loss of generality, we may assume they are integers with no common factor.
Then the side lengths $(X,Y,Z)$ are also integers, since they are integral linear
combinations of $(a,b,c)$.  
Let the sides of $T$ next to the $\alpha + \beta$ angle be $X$ and $Y$.  
Again we will show that $(X,Y,Z)$ are multiples of a triple with no common factor, and apply Lemma~\ref{lemma:claude3} 
and then the area equation.   Here are the details:

By the law of sines,
$$ X: {\sin 2\alpha} =   Y :{\sin 2\beta} = Z :{\sin (\alpha + \beta)}.$$
Using the formulas from Lemmas~\ref{lemma:claude} and~\ref{lemma:beta},  as well as $\alpha + \beta = \pi/3$, we have
$$ X: \frac {a \sqrt 3 (a+2b)}{2c^2} = Y:\frac {b \sqrt 3 (2a+b)}{2c^2} = Z: \frac{\sqrt 3 } 2 .$$
$$ X: a(a+2b) = Y : b(2a+b) = Z: c^2.$$
Then for some real $\lambda$ we have
\begin{equation}
 (X,Y,Z) = \lambda (a(a+2b), b(2a+b), c^2). \label{eq:618}
 \end{equation}
 We will show that $(a(a+2b), b(2a+b), c^2)$ have no common factor.  
 Suppose the prime $p$ divides all three. Then $p|c$, so $p \nmid a$ and $p \nmid b$,
 by Lemma~\ref{lemma:pairwisecoprimality}.
 Therefore $p \mid (2a + b)$.  Then
 \begin{eqnarray*}
 b \equiv -2a \pmod p   && \mbox{\qquad since $p \mid 2a+b$}\\
 0  \equiv 3a^2 \pmod p && \mbox{\qquad since $c^2 = a^2 + b^2 + ab$ and $c \equiv 0$}\\
 p = 3                   && \mbox{\qquad since $p \nmid a$} \\
 3 \nmid c               && \mbox{\qquad by  Lemma~\ref{lemma:not3divc}} \\
 p \nmid c               && \mbox{\qquad by the preceding two lines} 
 \end{eqnarray*}
 But that contradicts the assumption $p \mid c^2$.  That completes
 the proof that  $(a(a+2b), b(2a+b), c^2)$ have no common factor.
Then by Lemma~\ref{lemma:claude3}, $\lambda$ is a positive integer in (\ref{eq:618}).

The area equation  is 
$$ XY \sin (\alpha +\beta) \ = \ N ab \frac {\sqrt 3} 2.$$
Since $\alpha + \beta = \pi/3$, $\sin(\alpha+\beta) = \tfrac{\sqrt 3} 2$, so 
the area equation becomes
$$ XY = N ab.$$
Substituting the values for $X$ and $Y$ from (\ref{eq:618}) into the area equation,
we obtain
\begin{eqnarray*}
\lambda^2 a(a+2b) b(2a+b)&=& \ N ab  \\
N &=& \lambda^2 (a+2b) (2a+b)  
\end{eqnarray*}
Since $\lambda$ is an integer and $a$ and $b$ are positive, 
 $N$ is not prime.
\end{proof}

\section{The main theorem}


\begin{theorem} \label{theorem:main}
Let triangle $T$  be $N$-tiled by a tile not similar to $T$.  Suppose $N > 3$.
 Then $N$ is not prime.
\end{theorem}

\begin{proof}  If $T$ is isosceles, we are done by Theorem~\ref{theorem:isosceles}, part~(iii).
If $T$ is equilateral, we are done by Theorem~\ref{theorem:equilateral}.   If 
the angles of $T$ were commensurable, then by Theorem~\ref{theorem:l1995-5.3},
the tiling would be a reptiling.  But that contradicts the hypothesis that $T$ is 
not similar to the tile.   Therefore, the angles of $T$ are incommensurable.



By Theorem~\ref{theorem:l1995-4.1}, together with the hypothesis that $T$ is not similar to the tile, 
 the tiling must fall into Group~1 ($3\alpha + 2\beta =\pi$) or Group~2 ($\gamma = 2\pi/3$).
Also as discussed above, the Group~1 case has already been proved.  That leaves only tilings with 
$\gamma = 2\pi/3$ to consider.  According to Theorem~\ref{theorem:l1995-4.1}, $T$ must have
one of the following shapes:
$$(\alpha, 2\alpha, 3\beta), (\alpha, 2\beta, 2\alpha + \beta), (\alpha, \alpha+\beta, \alpha+2\beta), \text{or }  (2\alpha, 2\beta, \alpha+\beta).$$
In the preceding four theorems, it has been shown that in each of four cases, $N$ cannot be prime.
\end{proof}

\begin{corollary} \label{corollary:634}
Let $N$ be prime.  Then there is an $N$-tiling of some triangle if and only if $N = 2$, $N =3$, or 
$N \equiv 1 \pmod 4$.
\end{corollary}

\begin{proof} {\em Left to right}: Suppose $N$ is prime and there is an $N$-tiling of some triangle, and $N > 3$.
Then by Theorem~\ref{theorem:main},  the tiling is a reptiling. By Theorem~\ref{theorem:reptile},
$N$ is a square, or a sum of two squares, or three times a  square.  Since $N > 3$ and $N$ is prime,
it is not three times a square.  Then 
it must be a sum of two squares.  Then $N \equiv 1 \pmod 4$.
\smallskip

{\em Right to left}:  If $N = 2$, just divide any isosceles triangle in two by its altitude.  If $N=3$,
the $30-60-90$ triangle can be 3-tiled, as illustrated in Figure~\ref{figure:tiling13}.  If $N > 3$ is congruent to 1 mod 4,  it is a sum of two squares,
as is shown in every number theory textbook.  Let $N = e^2 + f^2$, with integers $e,f > 0$. 
Then the right triangle with legs $e,f$ can be reptiled into $N$ tiles, as illustrated for $N=13$ 
in Figure~\ref{figure:tiling3}.  In general, the altitude splits the triangle into two similar copies, subdivided quadratically into 
$e^2$ and $f^2$ tiles.
\end{proof}

\begin{remark} 
See Figure~\ref{figure:tiling13} for pictures of the exceptional prime tilings.
\end{remark}

\enlargethispage{3\baselineskip}





\section*{Use of AI}
We used Claude (Anthropic) for proof-checking and copy-editing. 
But also Lemma~\ref{lemma:claude3} was proved on
July 22, 2026 by the LLM, Claude Fable.  It was invented by Claude
in the course of answering the question, whether a triangle 
with angles $(\alpha,2\alpha,3\beta)$ with $\alpha$ not a rational
multiple of $\pi$ can be cut into a prime number of congruent triangles.
This prompt was given to Claude 
by Grigore Roșu.    The proof of Lemma~\ref{lemma:claude3}  presented here
is a hand-rewritten version of Claude's.  
This lemma is short, but important.
Also, Lemma~\ref{lemma:not3divc} was discovered and proved by Claude Fable.
The formulas for $N$ in the four last theorems were already in \cite{zhang2026},
but were rediscovered by Claude without looking at  \cite{zhang2026}.  

 \section*{Acknowledgement}
 Thanks to  Grigore Roșu for applying Claude Fable to this problem. 

% following helps with keeping plainurl from
% inserting a lot of ugly spaces.
\begingroup
\raggedright
\bibliographystyle{plainurl}
%\bibliography{TriangleTiling}
\begin{thebibliography}{10}

\bibitem{beeson-isosceles}
Michael Beeson.
\newblock Tilings of an isosceles triangle.
\newblock {\em Integers}.
\newblock to appear; \href{https://arxiv.org/abs/1206.1974}{arXiv:1206.1974
  [math.MG]}.

\bibitem{beeson-noseven}
Michael Beeson.
\newblock No triangle can be cut into seven congruent triangles.
\newblock url={https://arxiv.org/abs/1811.09723}, 2019.
\newblock \href {https://arxiv.org/abs/1811.09723} {\path{arXiv:1811.09723}}.

\bibitem{beeson-triangletiling3}
Michael Beeson.
\newblock Triangle tiling: The case $3\alpha + 2\beta = \pi$, 2019.
\newblock URL: \url{https://arxiv.org/abs/1206.2229}, \href
  {https://arxiv.org/abs/1206.2229} {\path{arXiv:1206.2229}}.

\bibitem{beeson-equilateral}
Michael Beeson.
\newblock Tiling an equilateral triangle.
\newblock url={https://arxiv.org/abs/1812.07014}, 2024.
\newblock \href {https://arxiv.org/abs/1812.07014} {\path{arXiv:1812.07014}}.

\bibitem{beeson-isoscelesPrime}
Michael Beeson.
\newblock No prime tiling of an isosceles triangle.
\newblock url={https://arxiv.org/abs/2607.19572}, 2026.
\newblock \href {https://arxiv.org/abs/2607.19572} {\path{arXiv:2607.19572}}.

\bibitem{BLZ2026}
Michael Beeson, Mikl{\'o}s Laczkovich, and Yan~X. Zhang.
\newblock Solution of {E}rd{\H{o}}s problem 633.
\newblock url={https://arxiv.org/abs/2604.03609}, 2026.
\newblock \href {https://arxiv.org/abs/2604.03609} {\path{arXiv:2604.03609}}.

\bibitem{beeson-zhang2026}
Michael Beeson and Yan~X Zhang.
\newblock Rationality of certain triangle tilings.
\newblock url={https://arxiv.org/abs/2604.01314}, 2026.
\newblock \href {https://arxiv.org/abs/2604.01314} {\path{arXiv:2604.01314}}.

\bibitem{erdos-633}
T.F. Bloom.
\newblock {Erdős Problem \#633}.
\newblock \url{https://www.erdosproblems.com/633},
\newblock accessed 2026-03-10.

\bibitem{erdos-634}
T.F. Bloom.
\newblock {Erdős Problem \#634}.
\newblock \url{https://www.erdosproblems.com/634},
\newblock accessed 2026-04-03.

\bibitem{laczkovich1995}
M.~Laczkovich.
\newblock {Tilings of triangles}.
\newblock {\em Discrete Mathematics}, 140:79--94, 1995.

\bibitem{laczkovich2012}
Mikl{\'o}s Laczkovich.
\newblock Tilings of convex polygons with congruent triangles.
\newblock {\em Discrete and Computational Geometry}, 38:330--372, 2012.
\newblock \href {https://doi.org/10.1007/s00454-012-9404-x}
  {\path{doi:10.1007/s00454-012-9404-x}}.

\bibitem{snover1991}
Stephen~L. Snover, Charles Waiveris, and John~K. Williams.
\newblock {Rep-tiling for triangles}.
\newblock {\em Discrete Mathematics}, 91:193--200, 1991.

\bibitem{soifer}
Alexander Soifer.
\newblock {\em How Does One Cut a Triangle?}
\newblock Springer, 2009.

\bibitem{zhang2026}
Yan~X Zhang.
\newblock Tiling triangles with $2\pi/3$ angles.
\newblock url={https://arxiv.org/abs/2512.22696}, 2026.
\newblock \href {https://arxiv.org/abs/2512.22696} {\path{arXiv:2512.22696}}.

\end{thebibliography}


\end{document}
